% !TEX root = ../main.tex
\lecture{9}{2026-11-12}{Produktion II}


\subsection{Langfristige Produktion}
Alle Faktoren sind variabel, inkl Kapital. Analog zu Nutzensfunktion ist die Langfristige Produktionsmengefunktion F eine multi-dimensionale Funktion

\subsubsection{Isoquanten}
Bei einer fixen Produktionsmenge ergeben alle Kombinationen von Kapital und Arbeit auf den Isoquanten (gleiche Menge) diesselbe Output Menge aus. \footnote{Es wäre die wissenschaftliche Aufgabe einer seriösen wirtschaftlichen Theorie der kap. Produktion, 
überhaupt zu erklären wieso Arbeit, die hier schon als abstrakte Kategorie als Arbeit allgemein, in einem einseitigen qualitativ gleichem Mass gemessen werden kann und wird. }

\subsubsection{Perfekte Input Komplemente}
Analog zur Nachfrage ergeben sich bei einseitiger Steigerung perfekter Kapital und Arbeit Komplemente keine grössere Produktionsmengen. Beide müssen stets in einem bestimmten Verhältnis erhöht werden, um eine grössere produzierte Menge erreichen zu können.
Menschliche Arbeit kann nicht mit Automatismen, Logistik, Material substituiert werden, um mehr produzieren zu können.
Hier ergeben sich konstante Skalenerträge, in dem man beide Elemente erhöht wird eine grössere Produktionsmenge ermöglicht


\subsubsection{Perfekte Input Substitute}
Im Beispiel der Arbeit könnten individuelle Produktivitätsunterschiede zwischen Arbeiter A und B sich in einem bestimmten Verhältnis austauschen lassen. Lineare Funktionen.

\subsubsection{Grenzrate der technischen Substitution (GRTS)}
Wie lassen sich bei gleicher Produktionsmenge die Faktoren gegenseitig substituieren, sodass die Menge unverändert bleibt bei jeder gegebner Faktor-Kombination
\begin{itemize}
    \item 
\end{itemize}


\textbf{Eigenschaften}

\begin{itemize}
    \item Konvexität
    \begin{itemize}
        \item Intuition: Je weniger ein Faktor relativ zum anderen schon eingesetzt ist, deto schwieriger ist es ihn noch mehr zu ersetzen.
        \item Beispiel: Maschinen können Arbeit leicht in einfachen Aufgaben ersetzen, die oft mehrheitlich sind, aber die wenige komplexe Aufgaben lassen sich nur gegen sehr grosse mengen an Kapital ersetzen, wobei menschliche Arbeit immer noch benötigt wird.
    \end{itemize}
\end{itemize}

\subsubsection{Skalenerträge}
\begin{itemize}
    \item zunehmend wenn:
    \begin{equation}
        F(zL,zK) > zF(L,K) | z^{\alpha + \beta} > z
    \end{equation}
    für alle (L,K) und z > 1. Also wenn eine vervielfachung der Faktoren in einer noch grösseren vervielfachung der Produktionsmenge verursacht
    \item konstante Skalenerträge
    \begin{equation}
        F(zL, zK) = zF(L,K) | z^{\alpha + \beta} = z \rightarrow \alpha = \beta
    \end{equation}
    Vervielfachung der Inputs resultiert in gleiche Vervielfachung des Outputs
    \item abnehmende Skalenerträge
    \begin{equation}
        F(zL,zK) < zF(L,K)
    \end{equation}
    Eine Vervielfachung des Inputs resultiert in reduzierte Outputs
\end{itemize}


\subsection{Langfristige Kosten}

